\begin{convention}[Index suppression and aggregation from the root variable]
\label{conv:notation}

Every variable in the model is defined with a fixed, ordered tuple of
\textbf{root indices} $(a_1, a_2, \dots, a_k)$, each ranging over a
given finite set; we call this the \textbf{root form}
$\phi_{a_1 a_2 \cdots a_k}$.
Writing a variable with fewer indices always denotes a
\emph{derived form}: one or more indices have been suppressed and the
root values aggregated according to the rules below.

Figure~\ref{fig:index-lattice} illustrates the full derivation lattice.

% ------------------------------------------------------------------
\medskip
\noindent\textbf{Rule~1 (Sum aggregation).
  Applies to all variable types: continuous, integer, and binary.}

Dropping index $a_\ell$ from any variable $\phi$ denotes the sum of
the root form over every value of that index:
%
\begin{equation*}
  \phi_{a_1 \cdots \widehat{a_\ell} \cdots a_k}
  \;\coloneqq\;
  \sum_{a_\ell} \phi_{a_1 \cdots a_\ell \cdots a_k},
\end{equation*}
%
where the hat $(\,\widehat{\cdot}\,)$ marks the suppressed position.
Multiple indices may be dropped simultaneously; each suppressed index
contributes one further summation.
For continuous and integer variables this is the natural interpretation.
For binary variables the result is an \emph{integer count} of how many
root-level binaries in the group equal~1; it is not itself binary.
No new variable is introduced by Rule~1, and no constraint is required.

% ------------------------------------------------------------------
\medskip
\noindent\textbf{Rule~2 (Dot notation for coincident index pairs).}

When two active root indices range over the \emph{same} set, suppressing
one while keeping the other would be ambiguous.
A dot~$(\,\cdot\,)$ marks the suppressed position explicitly:
%
\begin{equation*}
  \phi_{\cdot\, a_2 \cdots a_k}
  \;\coloneqq\;
  \sum_{a_1} \phi_{a_1 a_2 \cdots a_k},
  \qquad
  \phi_{a_1 \cdot\, a_3 \cdots a_k}
  \;\coloneqq\;
  \sum_{a_2} \phi_{a_1 a_2 \cdots a_k}.
\end{equation*}
%
The dot is only ever needed for indices belonging to the same set.
Rule~1 applies as normal to the result.

% ------------------------------------------------------------------
\medskip
\noindent\textbf{Rule~3 (OR-aggregated binaries).
  A new binary variable derived from the sum aggregate.}

When a binary indicator is needed at a coarser granularity---one that
signals whether \emph{at least one} root-level binary in the group is
active---a new binary decision variable is introduced, denoted with an
overline ($\overline{\phi}$).
This variable is not a notational shorthand: it must be declared in the
solver and its relationship to the root-level binaries enforced by an
explicit big-$M$ pair.
The pair is written directly in terms of the sum aggregate defined by
Rule~1, so that shorthand appears in the constraints themselves.

\smallskip
Suppressing index $a_\ell$ from a binary variable $\phi$ yields the
sum aggregate
$\phi_{a_1 \cdots \widehat{a_\ell} \cdots a_k}
 = \sum_{a_\ell} \phi_{a_1 \cdots a_\ell \cdots a_k}$
and the new binary $\overline{\phi}_{a_1 \cdots \widehat{a_\ell} \cdots a_k}$,
linked by:
%
\begin{align}
  \overline{\phi}_{a_1 \cdots \widehat{a_\ell} \cdots a_k}
    &\leq \phi_{a_1 \cdots \widehat{a_\ell} \cdots a_k}
  \label{eq:or-lb}\\
  \phi_{a_1 \cdots \widehat{a_\ell} \cdots a_k}
    &\leq M \cdot \overline{\phi}_{a_1 \cdots \widehat{a_\ell} \cdots a_k}
  \label{eq:or-ub}
\end{align}
%
Constraint~\eqref{eq:or-lb} forces $\overline{\phi}$ to~0 when the
count is~0 (no element active).
Constraint~\eqref{eq:or-ub} forces $\overline{\phi}$ to~1 when the
count is positive (at least one element active).
Together they enforce
$\overline{\phi} = 1 \iff \phi_{a_1 \cdots \widehat{a_\ell} \cdots a_k} \geq 1$.

% ------------------------------------------------------------------
\medskip
\noindent\textbf{Rule~4 (Exactly-one variant).
  An OR aggregate with an additional global restriction.}

In some cases the OR aggregates formed by Rule~3 must satisfy the
stronger condition that \emph{exactly one} value of a remaining index
$a_p$ is active for every fixed combination of the other indices.
Rather than introducing a further variable, we reuse the OR aggregate
$\overline{\phi}$ from Rule~3 and mark it with a hat ($\hat{\phi}$) to
signal that the following equality is imposed in addition to
constraints~\eqref{eq:or-lb}--\eqref{eq:or-ub}:
%
\begin{align}
  \sum_{a_p} \hat{\phi}_{a_1 \cdots \widehat{a_\ell} \cdots a_k}
    &= 1
  \label{eq:exactly-one}
\end{align}
%
That is, $\hat{\phi}$ and $\overline{\phi}$ are the \emph{same
variable}---the OR aggregate of $\phi$ over $a_\ell$, constructed
identically via Rules~1 and~3.
The hat is a diacritic warning the reader that
constraint~\eqref{eq:exactly-one} further restricts the feasible
pattern of values across $a_p$: whereas plain OR aggregation allows
any subset to be simultaneously active, the hat enforces that the sum
of those OR aggregates across all admissible values of $a_p$ equals
exactly one.
